\unitopen{13}
\begin{candos}
\item zapytać o drogę i zrozumieć odpowiedź (\de{geradeaus, links, rechts, an der Ampel});
\item podać drogę komuś innemu;
\item użyć przyimków z celownikiem (\de{mit, nach, zu, von, aus, bei, seit}) i zapamiętać \emph{dwa} z biernikiem (\de{durch, um});
\item kupić bilet i zapytać o połączenie.
\end{candos}

\section{Słownictwo}
\vocabtable{13}

\section{Przyimki, po których \emph{zawsze} jest celownik}
\de{aus, bei, mit, nach, seit, von, zu, gegenüber} wymagają \emph{celownika}, bez względu na to, czy pytasz \de{Wo?} czy \de{Wohin?}.
Pomaga rymowanka: „aus bei mit nach seit von zu”. Zlewanie: \de{zu dem = zum}, \de{zu der = zur}, \de{bei dem = beim},
\de{von dem = vom}.

\begin{center}\small
\begin{tabular}{@{}lll@{}}
\toprule
\textsf{\textbf{Przyimek}} & \textsf{\textbf{Użycie}} & \textsf{\textbf{Przykład}}\\
\midrule
\de{mit} & środek transportu, towarzystwo & \de{Ich fahre mit dem Bus.} \de{Ich komme mit meiner Kollegin.}\\
\de{nach} & miasto, kraj, dom; „po” & \de{nach Zürich}, \de{nach Hause}, \de{nach der Arbeit}\\
\de{zu} & osoba lub instytucja; cel & \de{zum Bahnhof}, \de{zur Arbeit}, \de{zu mir}\\
\de{von \dots{} bis} & zakres & \de{von Montag bis Freitag}\\
\de{aus} & pochodzenie, wnętrze & \de{aus Polen}, \de{aus dem Haus}\\
\de{bei} & przy kimś, w firmie & \de{bei der Firma Alpha}, \de{beim Arzt}\\
\de{seit} & od momentu (trwa nadal) & \de{seit einem Jahr}, \de{seit Montag}\\
\bottomrule
\end{tabular}
\end{center}
\textbf{Z biernikiem} (\de{durch}, \de{um}, \de{für}, \de{gegen}, \de{ohne}): \de{Ich gehe durch den Park.} \de{Ich fahre um die Ecke.}

\section{Droga}
\de{Wie komme ich zum Bahnhof?} Odpowiedzi: \de{Gehen Sie geradeaus.} \de{Biegen Sie an der Ampel links ab.}
\de{Nehmen Sie die zweite Straße rechts.} \de{Der Bahnhof ist gegenüber der Post.} \de{Es ist nicht weit, nur fünf Minuten zu Fuß.}
Transport: \de{Ich fahre mit dem Bus / dem Zug / dem Fahrrad}, ale \de{Ich gehe zu Fuß}.

\begin{falle}
\de{Ich fahre \textbf{zu} Zürich} jest błędem: miasta, kraje i dom wymagają \de{nach} (\de{nach Zürich, nach Hause}),
a \emph{osoby i instytucje} \de{zu} (\de{zum Arzt, zu Anna}). Drugi błąd: \de{seit} z dłuższym trwaniem to
\emph{teraźniejszość}, nie przeszłość: \de{Ich wohne seit einem Jahr in Zürich} (mieszkam od roku; po polsku też
teraźniejszość, ale uważaj na polskie „od roku mieszkałem”).
\end{falle}

\begin{szwajcaria}
Szwajcarska komunikacja publiczna (SBB, \de{ZVV} w Zurychu) jest punktualna i gęsta. Bilet to \de{das Billett}
(\de{die Fahrkarte} w DE). Rower to \de{das Velo}, tramwaj to \de{das Tram}, a \de{Halbtax} (półbilet) i \de{GA}
(\de{Generalabonnement}) to karty zniżkowe \verify{(warunki i ceny zmieniają się co roku, sprawdź u SBB)}.
Gdy pociąg się spóźnia, rozmowa jest krótka: \de{Der Zug hat Verspätung.}
\end{szwajcaria}

\begin{praca}
Dojazd do biura: \de{Ich fahre mit der S-Bahn zur Arbeit.} \de{Wie lange fahren Sie ins Büro?}
\de{Ich komme ein paar Minuten später.} Rozmowa o tym, jak do Ciebie dotrzeć, to pierwszy test praktyczny
z celownikiem: \de{Nehmen Sie die S-Bahn bis zum Hauptbahnhof.}
\end{praca}

\section{Dialogi}
\textbf{Dialog 1 --- pytanie o drogę}
\begin{dialog}
\turn{A}{Entschuldigung, wie komme ich zum Hauptbahnhof?}{Przepraszam, jak dojdę na dworzec główny?}
\turn{B}{Gehen Sie geradeaus bis zur Ampel. Dort biegen Sie links ab.}{Proszę iść prosto do świateł. Tam skręcić w lewo.}
\turn{A}{Ist es weit von hier?}{Czy to daleko stąd?}
\turn{B}{Nein, nur zehn Minuten zu Fuß. Oder nehmen Sie die Straßenbahn.}{Nie, tylko dziesięć minut pieszo. Albo proszę wziąć tramwaj.}
\turn{A}{Mit welcher Linie?}{Którą linią?}
\turn{B}{Mit der Linie 4. Die Haltestelle ist gegenüber der Post.}{Linią 4. Przystanek jest naprzeciwko poczty.}
\end{dialog}

\section{Ćwiczenia}
\begin{cwiczenia}
\exitem{Uzupełnij: \de{Ich fahre \dots{} dem Bus.}}{\de{mit}}
\exitem{Uzupełnij: \de{Wir fahren \dots{} Zürich.}}{\de{nach}}
\exitem{Uzupełnij: \de{Ich gehe \dots{} Arzt.} (zu + der Arzt, w skróconej formie)}{\de{zum}}
\exitem{Uzupełnij: \de{Ich komme \dots{} Polen.}}{\de{aus}}
\exitem{Uzupełnij: \de{Er arbeitet \dots{} einem Jahr bei Alpha.}}{\de{seit}}
\exitem{Uzupełnij: \de{Der Bahnhof ist \dots{} der Post.} (gegenüber)}{\de{gegenüber}}
\exitem{Uzupełnij: \de{Ich gehe \dots{} den Park.} (durch)}{\de{durch}}
\exitem{Popraw: \de{Ich fahre zu Zürich.}}{\de{Ich fahre nach Zürich.}}
\exitem{Popraw: \de{Ich gehe nach der Arzt.}}{\de{Ich gehe zum Arzt.}}
\exitem{Zapytaj o drogę do dworca.}{\de{Wie komme ich zum Bahnhof?}}
\exitem{Powiedz, że jedziesz do pracy tramwajem.}{\de{Ich fahre mit der Straßenbahn zur Arbeit.} (CH: \de{mit dem Tram})}
\exitem{Przetłumacz: „Skręć w lewo przy światłach.”}{\de{Biegen Sie an der Ampel links ab.} (nieformalnie: \de{Bieg an der Ampel links ab.})}
\end{cwiczenia}

\weekend{Opisz po niemiecku drogę z Twojego domu do pracy (lub do stacji, sklepu): minimum pięć zdań z
\de{mit}, \de{nach} lub \de{zu}, \de{geradeaus}, \de{links}, \de{rechts}. Potem poproś AI, aby zapytała o drogę, a Ty
odpowiedz bez notatek.}
